\section{总结与延伸}

\subsection{核心要点}

本讲从 VAE 出发，逐步推导出扩散模型的基本框架，核心要点包括：

\begin{enumerate}
    \item \textbf{VAE 的局限}：将似然和后验都建模为高斯过于简化，导致生成质量受限。
    \item \textbf{层级化思路}：引入多层潜变量，每层仍是简单高斯，但整体分布更复杂。
    \item \textbf{马尔可夫假设}：简化了层级化 VAE 的计算，使 ELBO 可以高效分解。
    \item \textbf{前向过程设计}：通过预定义的噪声调度 $\{\beta_t\}$ 逐步添加噪声，具有闭式解。
    \item \textbf{逆过程学习}：只需要学习去噪过程 $p_\theta(x_{t-1}|x_t)$，避免了联合训练的困难。
    \item \textbf{ELBO 分解}：先验匹配项自动满足，训练聚焦于一致性项。
\end{enumerate}

\subsection{拓展阅读}

\begin{itemize}
    \item Ho, J., Jain, A., \& Abbeel, P. (2020). \textit{Denoising Diffusion Probabilistic Models}. NeurIPS 2020. --- DDPM 原论文。
    \item Sohl-Dickstein, J., et al. (2015). \textit{Deep Unsupervised Learning using Nonequilibrium Thermodynamics}. ICML 2015. --- 扩散模型的理论基础论文。
    \item Kingma, D. P. \& Welling, M. (2014). \textit{Auto-Encoding Variational Bayes}. ICLR 2014. --- VAE 原论文。
    \item Luo, C. (2022). \textit{Understanding Diffusion Models: A Unified Perspective}. arXiv:2208.11970. --- 统一视角的综述。
    \item Nichol, A. Q. \& Dhariwal, P. (2021). \textit{Improved Denoising Diffusion Probabilistic Models}. ICML 2021. --- 关于余弦噪声调度等改进。
\end{itemize}

%% --- Fin del contenido --- %%

\end{document}
